\chapter{Cialis}

\section*{Definition and Overview}
Cialis is a brand name for tadalafil, a potent and selective phosphodiesterase type 5 (PDE5) inhibitor widely prescribed for the treatment of erectile dysfunction (ED) and the signs and symptoms of benign prostatic hyperplasia (BPH). By inhibiting the PDE5 enzyme, tadalafil prevents the degradation of cyclic guanosine monophosphate (cGMP) in the smooth muscle cells of the corpus cavernosum and prostate, promoting nitric oxide-mediated vasodilation. Cialis is unique among first-generation PDE5 inhibitors (such as sildenafil and vardenafil) due to its prolonged half-life of approximately 17.5 hours, which provides an extended therapeutic window of up to 36 hours. It is available in both on-demand and daily low-dose formulations, allowing for personalised management approaches.

\section*{Epidemiology}
Erectile dysfunction and benign prostatic hyperplasia are highly prevalent, age-dependent conditions affecting millions of men globally and in Australia. ED affects up to 20\% of men aged 40 to 59 and over 50\% of men aged over 70. BPH is similarly common, affecting approximately 50\% of men in their 50s and up to 90\% of men in their 80s. Cialis is one of the most frequently prescribed medications for these conditions in Australia. Due to the high co-occurrence of ED and lower urinary tract symptoms (LUTS) secondary to BPH in ageing men, Cialis offers a dual-action therapeutic option that simplifies treatment regimens and improves patient compliance.

\section*{Aetiology and Pathophysiology}
Erectile dysfunction and benign prostatic hyperplasia involve distinct yet overlapping physiological pathways that are receptive to PDE5 inhibition.
\begin{itemize}
    \item \textbf{Erectile dysfunction:} ED is caused by vascular, neurological, endocrine, or psychogenic factors. Vascular ED arises from endothelial dysfunction, atherosclerosis, or hypertension, which impair nitric oxide release and smooth muscle relaxation in the corpora cavernosa.
    \item \textbf{Benign prostatic hyperplasia:} Prostatic enlargement and bladder outlet obstruction lead to urinary hesitancy, weak stream, and nocturia. The pathophysiology involves increased smooth muscle tone in the bladder neck and prostate gland.
    \item \textbf{Role of phosphodiesterase type 5:} PDE5 is highly expressed in the penile corpus cavernosum, prostate, and bladder. Inhibiting PDE5 prevents the breakdown of cGMP, enhancing the vasodilatory effects of nitric oxide.
    \item \textbf{Vascular smooth muscle relaxation:} Accumulation of cGMP leads to intracellular calcium reduction and smooth muscle relaxation, facilitating blood inflow into the penis and reducing bladder neck tension.
    \item \textbf{Psychogenic and lifestyle influences:} Anxiety, depression, stress, smoking, alcohol misuse, and obesity can exacerbate ED by increasing sympathetic tone, which opposes nitric oxide-mediated vasodilation.
\end{itemize}

\section*{Clinical Features}
Cialis is prescribed to address specific clinical presentations of sexual dysfunction and lower urinary tract symptoms.
\begin{itemize}
    \item \textbf{Inability to attain or maintain an erection:} Difficulty achieving a penile erection sufficient for satisfactory sexual performance.
    \item \textbf{Lower urinary tract symptoms:} Symptoms including urinary frequency, urgency, hesitancy, weak urinary stream, dribbling, and nocturia associated with BPH.
    \item \textbf{Impact on quality of life:} Psychological distress, performance anxiety, relationship strain, and reduced self-esteem.
    \item \textbf{Physiological response window:} On-demand administration typically achieves efficacy within 30 to 60 minutes, lasting up to 36 hours in the presence of sexual stimulation.
\end{itemize}

\section*{Differential Diagnosis}
Before prescribing Cialis, it is essential to diagnose the primary cause of erectile dysfunction or urinary symptoms to ensure appropriate management.
\begin{itemize}
    \item \textbf{Cardiovascular disease:} ED is often an early sentinel marker of systemic atherosclerosis and coronary artery disease; a full cardiovascular assessment is required.
    \item \textbf{Endocrine disorders:} Hypogonadism (low testosterone), hyperprolactinaemia, and thyroid dysfunction can present with ED and must be ruled out or treated.
    \item \textbf{Neurogenic causes:} Conditions such as multiple sclerosis, Parkinson's disease, diabetic neuropathy, or history of pelvic surgery (e.g., radical prostatectomy) can impair neural pathways essential for erections.
    \item \textbf{Prostate cancer:} Must be differentiated from BPH via digital rectal examination and prostate-specific antigen (PSA) testing before initiating therapy.
    \item \textbf{Psychogenic erectile dysfunction:} ED stemming from anxiety, depression, or relationship issues rather than organic pathology.
\end{itemize}

\section*{When to Seek Medical Care}
Patients should consult a doctor to discuss treatment options for erectile dysfunction or urinary symptoms. Medical evaluation is crucial to screen for underlying cardiovascular risk factors.

\textbf{Call 000 (or local emergency services) for:}
\begin{itemize}
    \item A prolonged erection lasting more than 4 hours (priapism), which is a urological emergency that can lead to permanent tissue damage.
    \item Sudden, crushing chest pain (especially if Cialis has been taken recently; do NOT use nitroglycerin or other nitrates).
    \item Sudden decrease or loss of vision in one or both eyes.
    \item Sudden decrease or loss of hearing, sometimes accompanied by dizziness or tinnitus.
\end{itemize}

\section*{Diagnosis and Investigations}
A thorough diagnostic workup is required to confirm the indication and safety of Cialis.
\begin{itemize}
    \item \textbf{Clinical history and questionnaire:} Evaluation of symptoms using tools like the International Index of Erectile Function (IIEF) or the International Prostate Symptom Score (IPSS).
    \item \textbf{Cardiovascular assessment:} Blood pressure measurement, lipid panel, and evaluation of cardiac fitness for sexual activity.
    \item \textbf{Prostate-specific antigen test:} Performed to screen for prostate cancer in patients presenting with BPH symptoms.
    \item \textbf{Hormonal profile:} Total testosterone, prolactin, and thyroid-stimulating hormone (TSH) levels to rule out endocrine etiologies.
    \item \textbf{Digital rectal examination:} Conducted to assess the size, shape, and consistency of the prostate gland.
    \item \textbf{Medication reconciliation:} Crucial review to exclude patients taking nitrates, guanylate cyclase stimulators, or alpha-blockers, which can cause severe hypotension.
\end{itemize}

\section*{Management}
Cialis management involves personalised dosing and lifestyle optimization.
\begin{itemize}
    \item \textbf{On-demand dosing:} Tadalafil (typically 10\,mg or 20\,mg) taken at least 30 minutes before anticipated sexual activity.
    \item \textbf{Daily dosing:} Low-dose tadalafil (2.5\,mg or 5\,mg) taken once daily, suitable for frequent sexual activity or concurrent BPH management.
    \item \textbf{Lifestyle modifications:} Smoking cessation, regular aerobic exercise, weight reduction, and limiting alcohol intake to improve vascular health.
    \item \textbf{Combination therapies:} Concurrent use with non-pharmacological approaches, such as vacuum erection devices or psychosexual counseling.
    \item \textbf{Avoidance of drug interactions:} Strictly avoiding nitrates (e.g., glyceryl trinitrate, amyl nitrite) and cautious use with alpha-blockers (e.g., tamsulosin) to prevent severe blood pressure drops.
\end{itemize}

\section*{Complications}
Although generally well tolerated, Cialis can cause adverse effects and complications.
\begin{itemize}
    \item \textbf{Common side effects:} Headache, dyspepsia (indigestion), back pain, myalgia (muscle aches), nasal congestion, and facial flushing.
    \item \textbf{Severe hypotension:} Life-threatening blood pressure drops if co-administered with nitrates or guanylate cyclase stimulators (e.g., riociguat).
    \item \textbf{Priapism:} Prolonged, painful erection lasting over 4 hours, which can cause permanent erectile tissue damage.
    \item \textbf{Non-arteritic anterior ischaemic optic neuropathy:} A rare ophthalmic condition causing sudden vision loss, associated with PDE5 inhibitors.
    \item \textbf{Sudden sensorineural hearing loss:} Rare reports of sudden hearing reduction or loss.
\end{itemize}

\section*{Prognosis and Outcomes}
The prognosis for men treated with Cialis for ED and BPH is highly positive. The medication exhibits high efficacy, with over 70--80\% of men reporting significant improvements in erectile quality and satisfactory intercourse. In BPH, daily low-dose Cialis significantly reduces lower urinary tract symptoms within 2 to 4 weeks of starting therapy. Long-term compliance is high due to the flexible dosing options and favourable safety profile, provided patients are screened carefully for cardiovascular health and contraindications.

\section*{Prevention and Self-Care}
\begin{itemize}
    \item \textbf{Follow dosage instructions:} Take the medication exactly as prescribed; do not increase the dose without consulting your doctor.
    \item \textbf{Avoid recreational nitrates:} Never use "poppers" (amyl nitrite) while taking Cialis, as this combination can be fatal.
    \item \textbf{Limit alcohol consumption:} Excessive alcohol intake can temporarily impair erectile function and increase the risk of orthostatic hypotension.
    \item \textbf{Stay hydrated:} Drinking adequate water may help reduce the incidence of mild headaches or muscle aches.
    \item \textbf{Discuss all medications:} Inform your pharmacist or doctor about all over-the-counter and prescription medicines you are taking.
\end{itemize}

\section*{Sources and Further Reading}
The following reputable organisations provide further information on Cialis (tadalafil) and sexual health:
\begin{itemize}
    \item Healthdirect Australia. \textit{Consumer medicine information and guidelines for tadalafil.} https://www.healthdirect.gov.au/
    \item Healthy Male (Andrology Australia). \textit{Evidence-based resources on erectile dysfunction and reproductive health.} https://www.healthymale.org.au/
    \item Therapeutic Goods Administration (TGA). \textit{Australian prescription guidelines and safety alerts for PDE5 inhibitors.} https://www.tga.gov.au/
    \item Mayo Clinic. \textit{Detailed patient information on tadalafil administration and side effects.} https://www.mayoclinic.org/
    \item National Health Service (NHS). \textit{Guides on using tadalafil for erectile dysfunction and BPH.} https://www.nhs.uk/
\end{itemize}
